\begin{abstract}
Real-time speech systems answer most queries with a local model, but
some queries call for a stronger expert. Selective escalation is well
studied for text; speech tightens the constraint: the decision must be
made before the model commits to a response, and the signal must
survive distribution and modality shift. We ask whether a speech
model's own hidden states already carry such a competence signal. Our
key finding is that they do, but not where escalation methods usually
look: the conventional final-layer read collapses under distribution
shift, while a mid-network read stays reliable, transfers from
text-calibrated probes to speech input, and can be taken at a single
end-of-turn position before the first output token, in about
30\,ms. A linear probe on this representation gates a live escalation
loop: it raises deployed answer accuracy from 42\% to 63\% while
beating matched-rate random escalation, and, with weights, features,
and layer fixed, improves all five external speech-QA benchmarks over
always-local operation. These results identify a mid-network
representation, rather than the final-layer read, as the practical
signal for deciding when a speech model should hand off.
\end{abstract}
